\documentclass{article}

\PassOptionsToPackage{round,authoryear}{natbib}
\makeatletter
\def\input@path{{starter/}}
\makeatother
\usepackage[dblblindworkshop]{neurips_2026}

\usepackage[T1]{fontenc}
\usepackage{amsmath}
\usepackage{amsfonts}
\usepackage{booktabs}
\usepackage{graphicx}
\usepackage{microtype}
\usepackage{xcolor}
\usepackage{url}
\usepackage[hidelinks]{hyperref}

\graphicspath{{figures/}}
\newcommand{\model}{Qwen2.5-Math-1.5B-Instruct}
\newcommand{\dataset}{ProcessBench}

\setlength{\parskip}{0pt}
\setlength{\bibsep}{2.5pt}

\title{Decodability Is Not Localization:\
Auditing Hidden-State Error Probes in Qwen2.5-Math-1.5B-Instruct}
\workshoptitle{Interpretability for Discovery}
\author{Anonymous Author(s)}

\begin{document}
\maketitle

\begin{abstract}
We audit what a hidden-state error probe supports in a mathematical reasoning trace. In a case
study of \model{} on \dataset{}, a linear probe predicts the persistent target \emph{invalid so far}
with held-out AUROC $0.866$ [95\% CI $[0.849,0.884]$]. Equally tuned controls reach AUROC
$0.751$ with prefix TF-IDF, $0.776$ with structural metadata, and $0.810$ with both. A diagnostic
model that also receives final-answer correctness reaches $0.874$, but that field is a benchmark
annotation unavailable at inference time. A post-hoc nested comparison shows that adding the full
selected-layer hidden representation to the text and metadata model improves AUROC by $0.0587$
[95\% CI $[0.0414,0.0745]$], log loss by $-0.0798$ [95\% CI $[-0.1038,-0.0548]$], and Process F1 by
$+0.1866$ [95\% CI $[0.1397,0.2340]$] on the inspected split. The probe does not reliably localize the
first error: exact localization is $28.9\%$, correct-trace rejection is $61.2\%$, and detected
errors occur $0.59$ steps late on average. A post-hoc matched transition probe reaches AUROC
$0.769$ on 380 held-out pairs, but placebo reuse and covariate imbalance make this result
exploratory. The revised behavioral readout has AUROC $0.342$, so it fails the intervention
validity gate. The evidence supports incremental predictive information in this case study. It
does not establish generalization, calibration, self-monitoring, localization, or causal use.
\end{abstract}

\section{Introduction}

Process supervision evaluates the validity of intermediate steps rather than only the final answer
\citep{cobbe2021verifiers,lightman2023verify}. This creates a representational question: does a
model's hidden state change when a written solution has become invalid? A probe can answer a narrow
predictive version of that question. It cannot, by itself, show that the model uses the decoded
feature, nor that a threshold crossing identifies the first error. With a persistent label, every
boundary after an error is positive. A classifier can then use accumulated inconsistency, position,
solution length, generator identity, or final-answer outcome without responding at the onset.

We study these distinctions in one bounded setting. The model is \model{} and the data are 3,400
model-generated traces from \dataset{} \citep{yang2024qwenmath,zheng2025processbench}. The analysis
uses problem-grouped partitions, a persistent invalid-so-far target, hidden-state probes, visible
text and metadata controls, temporal randomization, matched transitions, a natural-token boundary
control, and a gated behavioral intervention. The analysis was designed as an audit of evidence
levels rather than as a claim that one score is a complete verifier.

The results have a specific shape. Hidden states predict the persistent target. They also add
information to a nuisance model on the inspected split when compared with the same model augmented
by the selected hidden representation. The score has temporal correspondence with the annotated
onset, but its threshold gives modest exact localization and many false alarms on correct traces.
The behavioral readout fails its AUROC gate, so the intervention is not interpreted as evidence of
causal use.

\section{Related work and contribution boundary}

Process reward models learn step quality from supervision on intermediate reasoning
\citep{lightman2023verify,li2023stepaware}. \dataset{} provides annotations of the earliest error
in mathematical reasoning traces and draws from GSM8K, MATH, OlympiadBench, and Omni-MATH
\citep{zheng2025processbench,hendrycks2021math,he2024olympiadbench,gao2024omnimath}. We use these
sources for one representation analysis. They are not independent datasets, and the analysis does
not test model-family generalization.

Several recent studies use internal states to score reasoning. Sun et al. probe arithmetic errors
and use the result for selective re-prompting \citep{sun2025arithmetic}. ReProbe studies hidden-state
step scoring for test-time scaling, and STEP uses a hidden-state scorer to prune reasoning traces
\citep{ni2026reprobe,liang2026earlysignals}. Alvarez and Baheri study first-error detection with
hidden-state transport geometry, including cross-model and cross-dataset evaluations
\citep{alvarez2026reasoningbreak}. Yuan et al. separate diagnostic error awareness from causal use
in a multi-model study \citep{yuan2026hiddenerror}. Our contribution is narrower. We report an
audit of one model and one benchmark construction that separates pooled prediction, first crossing,
temporal onset correspondence, source transfer, conditional prediction, and a behavioral validity
gate.

Probing literature distinguishes decodability from the information used by a model
\citep{hewitt2019control,belinkov2022probing,elazar2021amnesic}. Our intervention therefore has a
precondition: the unmodified behavioral readout must distinguish the relevant baseline states. The
paper reports the intervention records and the failed precondition without interpreting a dose
response as causal evidence.

\section{Experimental design}
\label{sec:method}

\subsection{Data, model, and partition}

Each \dataset{} trace contains a problem, generator identity, source, written steps, final-answer
correctness, and a human first-error index $e_i$. We set $e_i=-1$ for a fully correct trace. Traces
with the same normalized problem text are assigned to the same partition. Seed 42 creates train,
validation, and test partitions with target fractions of 60\%, 20\%, and 20\%. The completed data
flow retains 3,360 of 3,400 attempted traces and 24,909 of 25,697 attempted step boundaries. The
40 excluded traces exceed the 2,048-token context limit. The retained data contain 2,187 erroneous
traces and 1,173 fully correct traces. The held-out partition contains 669 traces, 432 erroneous
traces, 237 fully correct traces, and 4,985 step boundaries.

The retained test sources contain 84 GSM8K traces, 204 MATH traces, 188 OlympiadBench traces, and
193 Omni-MATH traces. The model record fixes \model{}, bfloat16 computation, a 2,048-token limit,
and the dataset fingerprint
\texttt{447f0a4b35c5747a9f9a3dab1e70d43}\allowbreak
\texttt{f71efd501497b8cec668b71337099784a}. The completed extraction used one NVIDIA A100 GPU
with batch size 16. The model reads each fixed trace in one causal forward pass and does not
produce the traces.

\subsection{Target and hidden-state probe}

At boundary $k$ of trace $i$, the primary target is
\begin{equation}
y_{ik}=\mathbb{1}[e_i\geq 0\ \wedge\ k\geq e_i].
\label{eq:target}
\end{equation}
This target means invalid so far. It labels the first invalid step and all later boundaries in an
erroneous trace. It does not identify whether the current step is locally incorrect.

The extraction records a 1,536-dimensional residual-stream state at each written-step boundary and
at 29 hidden-state indices. At each index, a standardized L2 logistic probe is trained with class
balancing. The regularization grid is $C\in\{0.01,0.1,1,10\}$. Validation AUROC selects $C$ and
validation Process F1 selects a threshold on a grid from 0.05 to 0.95 in increments of 0.005.
Validation selects index 23, $C=0.01$, and threshold $0.645$. The selected probe is refit on train
and validation data before the single test evaluation.

The predicted first error is the first threshold crossing, or $-1$ when no crossing occurs. We
report boundary AUROC, average precision, step F1, error exactness among erroneous traces,
correct-trace rejection, Process F1, complete-trace accuracy, detection rate, and error distance.
Primary probe intervals use 1,000 whole-trace bootstrap draws. The later CPU control audit uses
2,000 draws. The conditional hidden-state comparison uses 1,000 draws. These analyses are reported
as post-hoc audits rather than as one preregistered protocol.

\subsection{Visible controls and conditional comparisons}
\label{sec:controls}

The control analysis gives each nuisance model the same selection budget used for the probe. It
selects $C$ by validation AUROC, selects a threshold on validation data, refits on train and
validation data, and uses inverse boundary-count weights during fitting so each trace contributes
equal total training loss. Prefix TF-IDF uses the problem and reasoning text through the current
boundary. Structural metadata uses source, generator, absolute and fractional position, trace
length, and token count. The joint nuisance model combines prefix text and structural metadata.

Final-answer correctness is a benchmark annotation. It is available in the stored trace metadata,
but it is not an online signal for a detector operating before the final outcome is known. We use
models containing this field only as diagnostic upper bounds on nuisance predictability.

To test incremental prediction, we fit five feature sets. $N$ contains prefix TF-IDF and structural
metadata. $N+H$ adds the full selected-layer hidden representation. $O$ contains the nuisance
features and final-answer correctness. $O+H$ adds the full selected-layer hidden representation to
$O$. $H$ contains the hidden representation alone. Every feature set uses the same train,
validation, and test partitions, the same regularization grid, trace-equal training weights, and
refit on train plus validation. Paired intervals resample whole traces 1,000 times.

\subsection{Post-hoc temporal and onset analyses}

The temporal analyses were designed after inspecting the primary test result. We circularly shift
each frozen score trajectory 5,000 times, preserving its values while breaking its alignment to the
human annotation. We also compare the score jump at a noninitial first error with a transition in a
fully correct trace matched by source, generator, relative position, trace length, and token count.

The transition probe uses the hidden-state difference $h_{i,e_i}-h_{i,e_i-1}$. The fixed partitions
contain 1,152 train pairs, 387 validation pairs, and 380 test pairs. Across all partitions there are
1,919 pairs and 687 unique placebo traces. Placebo reuse has median 2, 90th percentile 6, and
maximum 46 pairs per trace. The standardized differences are 0.011 for relative position, 0.146
for trace length, and 0.285 for token count. Layer and $C$ are selected by validation AUROC, giving
index 21 and $C=0.01$. The test partition was inspected when this analysis was conceived, so the
result is exploratory.

A boundary-location control compares the last natural token of each written step with the last
token of its artificial step marker. It fixes index 23 and $C=0.01$, fits thresholds on validation
data, and evaluates on the original test partition. A paired whole-trace bootstrap compares the
natural-token and marker-token results.

\subsection{Behavioral validity gate}

The intervention moves the boundary state along the learned probe direction at layer 23 over
$\alpha\in\{-4,-2,-1,0,1,2,4\}$, where the scale is the training projection standard deviation.
It also evaluates 20 random directions orthogonal to the learned direction at the two extreme
values. The behavioral score is the difference between the conditional log probabilities of the
answers ``Yes'' and ``No'' for a question asking whether the reasoning contains an error.

The pipeline interprets a learned-direction effect only when the unmodified readout has AUROC
above 0.5 and nonzero specificity. The revised readout has specificity 0.203 but AUROC 0.342, so
it fails this gate. The intervention records are therefore reported as an assay diagnostic, not as
a causal test of model behavior.

\section{Results}

\subsection{Hidden states add information to a nuisance model}

The selected hidden probe reaches test AUROC $0.866$ [95\% CI $[0.849,0.884]$], average precision
$0.843$ [95\% CI $[0.813,0.870]$], and step F1 $0.743$ [95\% CI $[0.716,0.769]$]. It reaches error exactness
$0.289$, correct-trace rejection $0.612$, Process F1 $0.393$, and complete-trace accuracy
$0.404$.

The equally tuned nuisance models reach AUROC $0.751$ for prefix TF-IDF, $0.776$ for structural
metadata, and $0.810$ for their joint model. The outcome-augmented metadata model reaches $0.854$,
and the outcome-augmented joint model reaches $0.874$. The latter uses final-answer correctness,
so it is not a deployable baseline.

The conditional comparison gives the direct test of incremental prediction. $N$ alone reaches
AUROC $0.811$, log loss $0.533$, and Process F1 $0.211$. $N+H$ reaches AUROC $0.869$, log loss
$0.454$, and Process F1 $0.397$. The paired changes are AUROC $+0.0587$ [95\% CI
$[0.0414,0.0745]$], log loss $-0.0798$ [95\% CI $[-0.1038,-0.0548]$], and Process F1 $+0.1866$
[95\% CI $[0.1397,0.2340]$].

The oracle comparison gives $O$ AUROC $0.874$, log loss $0.440$, and Process F1 $0.284$. $O+H$
reaches AUROC $0.909$, log loss $0.385$, and Process F1 $0.444$. Its AUROC change is $+0.0348$
[95\% CI $[0.0242,0.0451]$]. The hidden AUROC advantage over the outcome-augmented joint nuisance
model is $-0.008$ [95\% CI $[-0.028,0.013]$], while the hidden probe has a Process F1 advantage of
$0.099$ [95\% CI $[0.040,0.158]$] and an error-exactness advantage of $0.113$ [95\% CI $[0.064,0.163]$].
The hidden probe has lower correct-trace rejection by $0.283$ [95\% CI $[-0.347,-0.218]$].

\begin{table}[t]
\centering
\small
\caption{Held-out predictive comparison. Nuisance controls receive the same validation selection
budget and trace-equal training weights. The final-answer field is a diagnostic annotation and is
not available to an online detector.}
\label{tab:controls}
\setlength{\tabcolsep}{3.5pt}
\begin{tabular}{lrrrr}
\toprule
Predictor & AUROC & Proc. F1 & Error exact & Corr. rej. \\
\midrule
Hidden state & .866 & .393 & .289 & .612 \\
Prefix TF-IDF & .751 & .274 & .192 & .477 \\
Structural metadata & .776 & .168 & .148 & .194 \\
Joint text + metadata & .810 & .210 & .169 & .278 \\
Metadata + final outcome & .854 & .262 & .164 & .646 \\
Joint + final outcome & .874 & .294 & .176 & .895 \\
\bottomrule
\end{tabular}
\end{table}

\begin{table}[t]
\centering
\small
\caption{Nested nuisance and hidden comparison. $\Delta$ is the hidden-augmented model minus its
nuisance-only counterpart. All models use the same selection budget and refit protocol.}
\label{tab:nested}
\setlength{\tabcolsep}{3.5pt}
\begin{tabular}{llrrrl}
\toprule
Nuisance set & Model & AUROC & Log loss & Proc. F1 & $\Delta$ AUROC \\
\midrule
$N$ & $N$ only & .811 & .533 & .211 & N/A \\
 & $N+H$ & .869 & .454 & .397 & $+$.0587 \\
$O$ & $O$ only & .874 & .440 & .284 & N/A \\
 & $O+H$ & .909 & .385 & .444 & $+$.0348 \\
\bottomrule
\end{tabular}
\\[-1mm]
\raggedright\footnotesize The paired 95\% AUROC intervals are $[0.0414,0.0745]$ for $N+H-N$
and $[0.0242,0.0451]$ for $O+H-O$. The $O$ model includes final-answer correctness.
\end{table}

Equal trace weighting at evaluation changes AUROC by $+0.0013$ [95\% CI
$[-0.0069,0.0103]$]. The corresponding changes are $-0.0052$ for average precision and
$-0.0087$ for step F1, with intervals that include zero. The length-aware threshold analysis
reduces Process F1 from $0.393$ to $0.377$ and correct-trace rejection from $0.612$ to $0.540$;
error exactness is unchanged at $0.289$.

\subsection{Localization is weaker than pooled discrimination}

The hidden score localizes $28.9\%$ of erroneous traces exactly, $59.3\%$ within one step, and
$72.7\%$ within two steps. It crosses on $88.7\%$ of erroneous traces and falsely crosses on
$38.8\%$ of fully correct traces. Among detected erroneous traces, the mean signed localization
error is $+0.59$ steps [95\% CI $[+0.36,+0.80]$]. The score therefore tends to identify a
persistent invalid state later than the annotated onset.

The temporal randomization test gives exact localization $0.289$ for the observed trajectories
and a shifted-null mean of $0.164$ [95\% null interval $[0.137,0.192]$], with permutation
$p=0.0002$. The observed within-one-step rate is $0.593$, compared with shifted-null mean $0.436$
and $p=0.0002$. The observed Process F1 is $0.393$, compared with shifted-null mean $0.259$ and
$p=0.0002$. These shifts break annotation alignment but do not remove position as an explanation.

The error-aligned score rises from $0.341$ immediately before the annotated error to $0.559$ at
the onset, $0.680$ one step later, and $0.731$ two steps later. The matched correct-trace
transition rises by $0.113$, while the error-onset transition rises by $0.257$. Their paired
difference is $0.144$ [95\% CI $[0.096,0.193]$]. Correct traces also show positive score drift.
The result is consistent with a gradual change near the annotation, not a sharp threshold event.

The natural-token boundary control gives AUROC $0.868$ and Process F1 $0.419$. The marker-token
control gives AUROC $0.866$ and Process F1 $0.404$. Natural-token minus marker-token differences
are $+0.002$ for AUROC [95\% CI $[-0.007,+0.011]$] and $+0.026$ for Process F1
[95\% CI $[-0.020,+0.075]$]. Both intervals include zero.

\subsection{Matched onset transitions are exploratory}

The transition probe reaches AUROC $0.769$ [95\% CI $[0.713,0.820]$], average precision $0.741$
[95\% CI $[0.672,0.816]$], and paired accuracy $0.753$ [95\% CI $[0.676,0.824]$] on 380 test pairs. The
position, destination-step TF-IDF, and shuffled-label controls reach AUROC $0.631$, $0.671$, and
$0.585$. The transition difference is therefore linearly distinguishable from the selected
controls on the inspected split.

The interpretation is limited by the matching diagnostics. One placebo trace supports as many as
46 pairs. The matched groups are close on relative position but differ more on trace length and
token count, with standardized differences 0.146 and 0.285. One-to-one matching and inverse-reuse
sensitivity refits are implemented but no completed sensitivity artifact is present. We therefore
treat this result as exploratory and do not use it to claim a general first-error mechanism.

\subsection{Ranking transfers across ProcessBench sources}

Source-specific probes trained on one source and evaluated on another have mean off-diagonal AUROC
$0.805$ and mean diagonal AUROC $0.845$. Every off-diagonal AUROC is above $0.739$. Thresholded
Process F1 is less stable, with mean off-diagonal value $0.261$, mean diagonal value $0.336$, and
off-diagonal values from $0.055$ to $0.371$. These are cross-source results within one benchmark
construction. They do not test another dataset, model family, or model size.

\subsection{Behavioral readout and intervention status}

The original single-token CORRECT versus INCORRECT readout has AUROC $0.283$ and specificity
zero. The revised conditional Yes versus No readout has AUROC $0.342$ and specificity $0.203$ on
256 balanced boundaries. It passes the specificity condition but fails the AUROC condition required
by the intervention pipeline. The learned-direction dose records therefore do not provide evidence
for or against causal use of the decoded feature. The intervention section is empty by design after
an invalid baseline readout, not because a valid behavioral test found a null effect.

\section{Scope and limitations}

The strongest supported statement is conditional: on one frozen split of one ProcessBench
construction, hidden states add held-out predictive information to a tuned nuisance model. This
statement is about prediction of the stored invalid-so-far label. It is not a claim that the model
has a deployable error detector, that the hidden representation is unavailable from visible text,
or that the model uses the feature in its own verdict.

The study uses one model, one benchmark construction, and one fixed partition. The post-hoc audits
were designed after the primary result. Their intervals describe resampling uncertainty within the
completed pipeline, not sensitivity to a new split, model, or selection procedure. The persistent
target labels all later boundaries after an annotated error. The visible controls do not cover every
possible textual or generator shortcut. The final-answer field is unavailable online. The source
transfer analysis measures ranking, while threshold-dependent metrics vary more strongly. Probability
calibration was not tested as a separate research question.

The transition analysis excludes 52 test traces whose first error is at step 0 and uses reused
placebo traces. The natural-token control and temporal randomization were also post-hoc. Counterfactual
activation patching has no result because drafted corrections were not independently verified and
its behavioral readout did not pass the validity gate.

The detector is not suitable for grading, certification, or unattended feedback. On fully correct
traces, the selected threshold produces a $38.8\%$ false-alarm rate. Exact localization among
erroneous traces is $28.9\%$, and the mean detected error is late. The stored results support use as
a research diagnostic paired with human or formal verification only.

\section{Conclusion}

In \model{} reading \dataset{} traces, invalid-so-far prefixes are linearly decodable. The nested
comparison shows an incremental held-out gain from adding the selected hidden representation to
prefix text and structural metadata on the inspected split. The gain is not a complete verifier:
localization remains weak, correct-trace rejection is limited, threshold behavior varies across
sources and lengths, and the native behavioral readout fails its AUROC gate. The appropriate claim
is therefore incremental predictive information in a bounded case study, not localization,
self-monitoring, generalization, calibration, or causal use.

\bibliographystyle{plainnat}
\bibliography{references}

\appendix

\section{Data accounting}
\label{app:data}

\begin{table}[h]
\centering
\small
\caption{Attempted and retained traces and boundaries by frozen partition.}
\label{tab:dataflow}
\begin{tabular}{lrrrrrr}
\toprule
Partition & Attempted & Retained & Excluded & Boundaries & Error traces & Correct traces \\
\midrule
Train & 2,041 & 2,014 & 27 & 14,944 & 1,311 & 703 \\
Validation & 681 & 677 & 4 & 4,980 & 444 & 233 \\
Test & 678 & 669 & 9 & 4,985 & 432 & 237 \\
\midrule
All & 3,400 & 3,360 & 40 & 24,909 & 2,187 & 1,173 \\
\bottomrule
\end{tabular}
\end{table}

\begin{table}[h]
\centering
\small
\caption{Attempted, retained, and excluded traces by source and frozen partition.}
\begin{tabular}{lrrrr}
\toprule
Partition & GSM8K & MATH & OlympiadBench & Omni-MATH \\
\midrule
Train & 236/236/0 & 597/592/5 & 602/594/8 & 606/592/14 \\
Validation & 80/80/0 & 198/198/0 & 204/202/2 & 199/197/2 \\
Test & 84/84/0 & 205/204/1 & 194/188/6 & 195/193/2 \\
\bottomrule
\end{tabular}
\end{table}

The retained boundary counts by source are 1,229, 3,718, 5,129, and 4,868 for train; 414, 1,303,
1,708, and 1,555 for validation; and 439, 1,387, 1,661, and 1,498 for test, in the source order
shown in the table. The original attempted total was 25,697 boundaries. The retained total is
24,909 after excluding 40 traces that exceed the context limit.

\section{Descriptive subgroup audit}
\label{app:subgroups}

\begin{table}[h]
\centering
\small
\caption{Trace-macro outcomes by ProcessBench source with 95\% whole-trace bootstrap intervals.}
\begin{tabular}{lrrr}
\toprule
Source & Complete accuracy & Exact error & Correct rejection \\
\midrule
GSM8K & .512 [.405,.619] & .238 [.116,.375] & .786 [.658,.902] \\
MATH & .500 [.426,.569] & .347 [.264,.435] & .723 [.620,.818] \\
OlympiadBench & .298 [.234,.362] & .207 [.138,.282] & .463 [.340,.585] \\
Omni-MATH & .358 [.290,.425] & .324 [.253,.401] & .467 [.319,.614] \\
\bottomrule
\end{tabular}
\end{table}

Among traces with an incorrect final answer, complete-trace accuracy is $0.292$ [95\% CI $[0.245,0.342]$].
Among traces with a correct final answer, it is $0.506$ [95\% CI $[0.451,0.554]$]. Across generators with
at least 20 held-out traces, complete-trace accuracy ranges from $0.316$ to $0.621$. These are
descriptive subgroup results, not generator rankings or independent hypothesis tests.

\section{Reproduction record}
\label{app:repro}

The maintained package uses one configuration and one Click command group. The principal commands
are:
\begin{verbatim}
uv sync --extra dev
uv run math-error --config configs/project.yaml validate-config
uv run math-error --config configs/project.yaml run-all
uv run math-error --config configs/project.yaml analyze
uv run math-error --config configs/project.yaml fit-conditional
uv run math-error --config configs/project.yaml fit-contextual-baseline
uv run math-error --config configs/project.yaml fit-transition
uv run math-error --config configs/project.yaml transition-diagnostics
uv run math-error --config configs/project.yaml transition-sensitivity
uv run pytest
uv run ruff check src tests
\end{verbatim}

The project configuration records seed 42, the model and dataset identities, the 2,048-token limit,
selection grids, bootstrap counts, and artifact paths. The extraction and analysis artifacts record
selected layers, thresholds, predictions, metric summaries, bootstrap intervals, and the data flow.
The current repository also includes the unified package layout and neutral test layout described in
\texttt{README.md}.

\begin{table}[h]
\centering
\small
\caption{Compute and software record.}
\label{tab:compute}
\begin{tabular}{p{0.26\linewidth}p{0.64\linewidth}}
\toprule
Activation extraction & One NVIDIA A100 GPU, bfloat16, batch size 16, single causal passes over the retained traces. \\
Probe fitting and audits & CPU-capable; the control audit uses 2,000 whole-trace bootstrap draws. \\
Software record & Python 3.10 or newer; dependencies are declared in \texttt{pyproject.toml} and resolved in \texttt{uv.lock}. \\
Runtime record & The CPU audit records Python, NumPy, pandas, SciPy, scikit-learn, matplotlib, and platform versions in its \texttt{summary.json}. \\
Identifiers & The extraction record fixes \model{}, bfloat16 dtype, the 2,048-token limit, and the dataset SHA-256. \\
Licensing & Dataset and model licensing follows the published terms of the corresponding sources. \\
\bottomrule
\end{tabular}
\end{table}

\newpage
\input{checklist.tex}

\end{document}
